\documentclass[a4paper,10pt]{article}
\usepackage[utf8]{inputenc}
\usepackage[T1]{fontenc}
\usepackage{lmodern}
\usepackage[ngerman]{babel}
\usepackage[a4paper,portrait,left=22mm,right=22mm,top=22mm,bottom=24mm]{geometry}
\usepackage{longtable,array,booktabs}
\usepackage{xcolor}
\usepackage{enumitem}
\usepackage{microtype}
\usepackage{fancyhdr}
\usepackage{amsmath}
\usepackage[font=small,skip=5pt]{caption}
\pagestyle{fancy}
\fancyhf{}
\fancyfoot[L]{\small\texttt{C3\_SPM\_Blatt5.pdf}}
\fancyfoot[R]{\small Seite \thepage}
\renewcommand{\headrulewidth}{0pt}
\renewcommand{\footrulewidth}{0.4pt}
\setlength{\footskip}{22pt}
\setlength{\parindent}{0pt}
\setlength{\parskip}{4pt}
\setlength{\emergencystretch}{3em}
\setlist{nosep,leftmargin=*,topsep=3pt}
\clubpenalty=10000
\widowpenalty=10000
\displaywidowpenalty=10000
\raggedbottom
\setlength{\tabcolsep}{4pt}
\renewcommand{\arraystretch}{1.15}
\setlength{\LTpre}{4pt}
\setlength{\LTpost}{8pt}
\setlength{\LTcapwidth}{\textwidth}
\newcommand{\sectionspace}[1]{\par\ifdim
                                      \dimexpr\pagegoal-\pagetotal\relax<#1\relax\newpage
\fi\addvspace{14pt}}
\newcommand{\offen}{\emph{noch offen}}


\title{Qualitätsmanagement, Abschlusspräsentation und Qualitätsbericht\\{\large Übungsblatt 5}}
\author{Gruppe C3 / MoinMetrics\\{}Alexander Borrmann\\{}Lukas Pancratius\\{}Ino Wölke}
\date{Hochschule Emden-Leer\\Software-Projektmanagement\\Wintersemester 2026/27\\Prof. Dr. Gerrit Jan Veltink\\[1em]Abgabedatum: noch offen}
\begin{document}
\maketitle
\thispagestyle{fancy}
\clearpage
\section*{Aufgabe 1: Projektvorlage zur Qualitätssicherung}
% Eigenes Produkt der zugeordneten Gruppe vorlegen; Termine in Moodle beachten.
\subsection*{Zuordnung und Übergabe}
\offen

\section*{Aufgabe 2: Qualitätssicherung}
% Das Produkt der zugeordneten Gruppe prüfen und alle entdeckten Mängel notieren.
% Softwarefehler, Unklarheiten und Schwierigkeiten bei der Anwendung erfassen.
% Reine Geschmacksfragen und abweichende Implementierungen sind keine Mängel.
\subsection*{Geprüftes Produkt und Testbedingungen}
\offen

\subsection*{Festgestellte Mängel}
\offen

\section*{Aufgabe 3: Abschlusspräsentation und Qualitätsbericht}
\subsection*{Abschlusspräsentation des eigenen Projekts}
% Kritische Auseinandersetzung mit den Themen des Softwareprojektmanagements.
\offen

\subsection*{Qualitätsbericht zum geprüften Produkt}
% Als Präsentation direkt nach der Projektpräsentation der zugeordneten Gruppe.
% Abgabe gleichzeitig mit der eigenen Abschlusspräsentation.
\offen

\subsection*{Finale Abgabe in Moodle}
% Separate Abgabebestandteile; dieses Gerüst ersetzt diese Dateien nicht.
\begin{enumerate}
    \item Funktionsfähiges Produkt als WAR-Datei: \offen
    \item Quellcode des Produkts: \offen
    \item Executive Abstract für die Marketingmanager*innen (maximal eine Seite): \offen
    \item Präsentation der Ergebnisse: \offen
    \item Bericht der Qualitätssicherung: \offen
\end{enumerate}
\end{document}
